% Lean module: MonadFormal.Temporal
\section{Temporal dependence and whole-path constraints}
\label{sec:temporal-elimination}

A future-facing formula need not describe a new causal law. The same
ordinary differential system can yield either a past-memory expression
or a terminal-value expression, according to the boundary information
used to eliminate an unobserved component. This gives a concrete null
explanation for one part of the whole-path intuition. Projection and
memory also appear in the broader Mori--Zwanzig programme
\cite{chorin2000optimal}; the result below is the bounded, linear case,
not the general nonlinear theory.

Three cases must be kept separate: elimination of an ordinary hidden
state, a literal mixed advance--delay equation, and a variational law
on complete paths. Each has its own data and existence question.

\subsection{Eliminating a hidden linear state}

Let $U,V$ be real Banach spaces and let
\[
 A:U\to U,\quad B:V\to U,\quad C:U\to V,\quad D:V\to V
\]
be bounded linear operators. Define the operator-exponential propagator
$E_D(t)=\exp(tD)$. The paths below are defined on $\mathbb R$ and are
differentiable at every point of the stated closed interval. This
endpoint assumption excludes merely one-sided endpoint derivatives.
All integrals are Bochner
integrals, with their usual orientation.
\leandef{MonadFormal.Temporal.flow}

\begin{theorem}[Variation of constants and terminal anchoring]
\label{thm:temporal-variation}
If $f$ is continuous on the interval between $a$ and $t$ and
$v'(s)=f(s)+Dv(s)$ throughout that interval, then
\begin{equation}\label{eq:temporal-variation}
 v(t)=E_D(t-a)v(a)+\int_a^t E_D(t-s)f(s)\,ds.
\end{equation}
In particular, on $[t,T]$ the same hypotheses give
\begin{equation}\label{eq:temporal-terminal}
 v(t)=E_D(t-T)v(T)-\int_t^T E_D(t-s)f(s)\,ds.
\end{equation}
\end{theorem}
\begin{proof}
The exponential is continuous, $E_D(0)=I$, and
$E_D'(r)=E_D(r)D$. The product and chain rules therefore give
\[
 \frac{d}{ds}\bigl(E_D(t-s)v(s)\bigr)
 =-E_D(t-s)Dv(s)+E_D(t-s)\bigl(f(s)+Dv(s)\bigr)
 =E_D(t-s)f(s).
\]
The last integrand is continuous and hence integrable. Integrating
between $a$ and $t$ gives \eqref{eq:temporal-variation}. Anchor that
identity at $T$ and reverse the oriented integral to obtain
\eqref{eq:temporal-terminal}.
\leanproof{MonadFormal.Temporal.flow_zero,MonadFormal.Temporal.continuous_flow,MonadFormal.Temporal.flow_derivative,MonadFormal.Temporal.backward_flow_derivative,MonadFormal.Temporal.weighted_forcing_integrable,MonadFormal.Temporal.integrating_factor_derivative,MonadFormal.Temporal.variation_of_constants,MonadFormal.Temporal.terminal_variation_of_constants}
\end{proof}

\begin{theorem}[Exact memory and future-integral reductions]
\label{thm:temporal-reductions}
Suppose $u$ is continuous and $v'=Cu+Dv$ on $[a,t]$, and
$u'(t)=Au(t)+Bv(t)$. Then
\begin{equation}\label{eq:temporal-memory}
 u'(t)=Au(t)+BE_D(t-a)v(a)
       +\int_a^t BE_D(t-s)Cu(s)\,ds.
\end{equation}
With the corresponding hypotheses on $[t,T]$, instead,
\begin{equation}\label{eq:temporal-future}
 u'(t)=Au(t)+BE_D(t-T)v(T)
       -\int_t^T BE_D(t-s)Cu(s)\,ds.
\end{equation}
If both component derivatives hold at $t$, the paired state satisfies
\[
 \frac{d}{dt}(u(t),v(t))
 =\bigl(Au(t)+Bv(t),\;Cu(t)+Dv(t)\bigr).
\]
\end{theorem}
\begin{proof}
Apply \Cref{thm:temporal-variation} with $f=Cu$, substitute into
$Au(t)+Bv(t)$, and commute the bounded linear map $B$ with the
integral. For the terminal form, use \eqref{eq:temporal-terminal}.
The derivative of the paired state is the pair of component derivatives.
\leanproof{MonadFormal.Temporal.memory_equation,MonadFormal.Temporal.terminal_equation,MonadFormal.Temporal.augmented_derivative}
\end{proof}

These are identities obeyed by a coupled-system solution. They do not
assert existence for arbitrary mixed boundary data, nor do they make
$v(a)$ or $v(T)$ accessible to an observer who did not retain or obtain
it. Enlarging the mathematical state to $(u,v)$ does not establish an
affordable encoding in the operational access state $A_t$.
Equation~\eqref{eq:temporal-future} is a terminally constrained
description of ordinary dynamics, not an initial-value update rule that
reads an independently chosen future. Literal advance--delay dynamics
and whole-path variational constraints require separate well-posedness
and boundary analyses.

% Lean module: MonadFormal.Temporal.MixedBoundary
\subsection{A literal mixed equation and a resonance obstruction}
\label{sec:mixed-boundary}

The eliminated-state identities do not settle well-posedness of an
arbitrary equation containing advanced arguments. Consider the scalar
boundary problem
\begin{equation}\label{eq:mixed-boundary}
 x'(t)=\alpha\bigl(x(t-\tfrac12)+x(t+\tfrac12)\bigr),\qquad
 x(t)=0\ (t\le0),\quad x(t)=1\ (t>1).
\end{equation}
Here a solution means a real-valued path continuous on $[0,1]$, with
the displayed derivative on the two open intervals $(0,\tfrac12)$ and
$(\tfrac12,1)$. We do not require differentiability at $\tfrac12$ or
continuity across the future boundary at $1$. These qualifications are
part of the problem, not conclusions of the calculation.
\leandef{MonadFormal.Temporal.MixedBoundary.IsSolution}

Put $u(s)=x(s)$ and $v(s)=x(s+\tfrac12)$ for $0\le s\le\tfrac12$.
The two open-interval equations become
\[
 u'=\alpha v,\qquad v'=\alpha(u+1),\qquad
 u(0)=0,\quad u(\tfrac12)=v(0).
\]
The last equality is the seam condition: both sides describe
$x(\tfrac12)$. Define
\begin{align*}
 u_\beta(s)&=\cosh(\alpha s)-1+\beta\sinh(\alpha s),\\
 v_\beta(s)&=\sinh(\alpha s)+\beta\cosh(\alpha s).
\end{align*}
\leandef{MonadFormal.Temporal.MixedBoundary.left,MonadFormal.Temporal.MixedBoundary.right}

\begin{lemma}[Exhaustive split-system solution]
\label[lemma]{lem:mixed-split}
Suppose $u,v$ are continuous on $[0,t]$, differentiable on $(0,t)$,
satisfy $u'=\alpha v$, $v'=\alpha(u+1)$ there, and $u(0)=0$.
For $t\ge0$, their endpoint values are
$u(t)=u_\beta(t)$ and $v(t)=v_\beta(t)$, where $\beta=v(0)$.
Conversely, $u_\beta,v_\beta$ satisfy the two differential equations
everywhere and have initial values $0,\beta$.
\end{lemma}
\begin{proof}
For $z'=cz$, the derivative of $z(s)e^{-cs}$ is zero on the
interior. Continuity at the endpoints and the fundamental theorem of
calculus give $z(t)=z(0)e^{ct}$. Apply this first to $u+1+v$, with
$c=\alpha$, and then to $u+1-v$, with $c=-\alpha$. Adding and
subtracting the resulting identities gives the displayed hyperbolic
formulas. Their derivatives and initial values follow directly from
the derivatives of $\sinh$ and $\cosh$.
\leanproof{MonadFormal.Temporal.MixedBoundary.scalar_solution,MonadFormal.Temporal.MixedBoundary.split_solution,MonadFormal.Temporal.MixedBoundary.left_zero,MonadFormal.Temporal.MixedBoundary.right_zero,MonadFormal.Temporal.MixedBoundary.left_derivative,MonadFormal.Temporal.MixedBoundary.right_derivative}
\end{proof}

\begin{theorem}[Boundary compatibility and exact existence criterion]
\label{thm:mixed-existence}
Write $S=\sinh(\alpha/2)$ and $C=\cosh(\alpha/2)$.
The seam condition is equivalent to
\[
 (1-S)\beta=C-1.
\]
If $S\ne1$, it selects exactly one coefficient,
$\beta=(C-1)/(1-S)$, and \eqref{eq:mixed-boundary} has a solution in
the declared class. If $S=1$, there is no solution in that class.
In particular, $\alpha=2\operatorname{arsinh}(1)$ is an explicit
resonant parameter at which existence fails.
\end{theorem}
\begin{proof}
Substitute the explicit formulas into $u_\beta(\tfrac12)=v_\beta(0)$.
When $S\ne1$, divide by $1-S$ and assemble
\[
 x_\beta(t)=
 \begin{cases}
 0,&t\le0,\\
 u_\beta(t),&0<t\le\tfrac12,\\
 v_\beta(t-\tfrac12),&\tfrac12<t\le1,\\
 1,&t>1.
 \end{cases}
\]
The initial and seam equalities make this path continuous on $[0,1]$.
On the first open half its delayed value is zero and its advanced
value is $v_\beta(t)$; on the second, those values are
$u_\beta(t-\tfrac12)$ and one. The split derivatives therefore give
\eqref{eq:mixed-boundary} on both open halves.

At $S=1$, the seam equation would require $C=1$, contradicting
$C^2-S^2=1$. This excludes every solution, not just the explicit
construction: any solution gives continuous split paths, and
\Cref{lem:mixed-split} forces their values at the seam to have the
same form. Finally, $\sinh(\operatorname{arsinh}(1))=1$ supplies the
stated parameter.
\leandef{MonadFormal.Temporal.MixedBoundary.path}
\leanproof{MonadFormal.Temporal.MixedBoundary.join_iff,MonadFormal.Temporal.MixedBoundary.join_unique,MonadFormal.Temporal.MixedBoundary.no_join_at_resonance,MonadFormal.Temporal.MixedBoundary.resonant_parameter,MonadFormal.Temporal.MixedBoundary.no_split_at_resonance,MonadFormal.Temporal.MixedBoundary.path_history,MonadFormal.Temporal.MixedBoundary.path_left,MonadFormal.Temporal.MixedBoundary.path_right,MonadFormal.Temporal.MixedBoundary.path_future,MonadFormal.Temporal.MixedBoundary.path_equation_left,MonadFormal.Temporal.MixedBoundary.path_equation_right,MonadFormal.Temporal.MixedBoundary.path_continuousOn,MonadFormal.Temporal.MixedBoundary.exists_mixed_solution,MonadFormal.Temporal.MixedBoundary.no_mixed_at_resonance,MonadFormal.Temporal.MixedBoundary.solution_exists_iff}
\end{proof}

\begin{corollary}[Matching the future endpoint is an extra constraint]
For the joined, non-resonant construction, $v_\beta(\tfrac12)=1$
if and only if $C=2S$, equivalently $\alpha=\log3$.
\end{corollary}
\begin{proof}
Multiply $v_\beta(\tfrac12)=S+\beta C$ by $1-S$, use the seam
equation and $C^2-S^2=1$, and obtain
\[
 v_\beta(\tfrac12)(1-S)=1+S-C.
\]
Since $1-S\ne0$, the desired value is equivalent to $C=2S$.
Writing the hyperbolic functions as exponentials then gives
$e^\alpha=3$, equivalently $\alpha=\log3$.
\leanproof{MonadFormal.Temporal.MixedBoundary.terminal_match_iff,MonadFormal.Temporal.MixedBoundary.terminal_parameter_iff}
\end{proof}

Thus the solution convention matters. A path in the declared class may
jump immediately after $1$; requiring a continuous join to the unit
future imposes the additional endpoint equation. Neither the explicit
construction nor its resonance obstruction establishes a general
well-posedness theorem for advanced functional differential equations.

% Lean module: MonadFormal.Temporal.Variational
\subsection{A coercive whole-path energy}
\label{sec:temporal-variational}

Whole-path consistency is a third interpretation of temporal dependence.
It specifies an extremum subject to boundary data, rather than an
initial-value update rule. The following Hilbert-space result supplies
the existence argument for the interval construction below.

Let $V$ be a real Hilbert space, $W$ a real inner-product space,
$T:V\to W$ a bounded linear operator, and $b\in W$. Put
\[
 E_\lambda(u)=\frac12\lVert u\rVert^2
       +\frac\lambda2\lVert Tu+b\rVert^2,\qquad
 B_\lambda(u,h)=\langle u,h\rangle+\lambda\langle Tu,Th\rangle.
\]
\leandef{MonadFormal.Temporal.Variational.energy,MonadFormal.Temporal.Variational.form}

\begin{theorem}[Coercive quadratic minimization]
\label{thm:variational-minimum}
If $\lambda\ge0$, then $E_\lambda$ has a unique global minimizer.
There is a stationary point $u_*$ satisfying
\begin{equation}\label{eq:variational-weak}
 \langle u_*,h\rangle+\lambda\langle Tu_*+b,Th\rangle=0
 \quad\text{for every }h\in V.
\end{equation}
Every such point obeys the exact gap identity
\begin{equation}\label{eq:variational-gap}
 E_\lambda(u_*+h)-E_\lambda(u_*)
 =\frac12\bigl(\lVert h\rVert^2+\lambda\lVert Th\rVert^2\bigr).
\end{equation}
\end{theorem}
\begin{proof}
The bounded bilinear form satisfies
$B_\lambda(u,u)\ge\lVert u\rVert^2$.
The functional $h\mapsto-\lambda\langle b,Th\rangle$ is bounded,
so Lax--Milgram gives \eqref{eq:variational-weak}.
Expand the two squared norms at $u_*+h$; the linear cross terms
vanish by stationarity, giving \eqref{eq:variational-gap}.
Its right-hand side is nonnegative and is zero only if $h=0$.
Thus $u_*$ minimizes the energy, and any other minimizer equals it.
\leanproof{MonadFormal.Temporal.Variational.form_apply,MonadFormal.Temporal.Variational.form_coercive,MonadFormal.Temporal.Variational.exists_stationary,MonadFormal.Temporal.Variational.energy_gap,MonadFormal.Temporal.Variational.minimizes_of_stationary,MonadFormal.Temporal.Variational.exists_unique_minimizer}
\end{proof}

% Lean module: MonadFormal.Temporal.Variational.Interval
\subsubsection{Endpoint-constrained interval paths}
\label{sec:variational-interval}

Let $\mu$ be Lebesgue measure restricted to $[0,1]$, and set
$V=L^2(\mu;\mathbb R)$. For $f\in V$, define
\[
 (Pf)(t)=\int_0^t f(s)\,d\mu(s),\qquad
 H=\{f\in V:(Pf)(1)=0\}.
\]
Functions in $V$ are identified almost everywhere. Their primitives
are ordinary continuous functions; the integral does not depend on
the representative. Given endpoint values $a,b\in\mathbb R$, put
\[
 x_f(t)=a+(b-a)t+(Pf)(t),\qquad f\in H.
\]
The restricted measure extends the primitive constantly outside the
unit interval; this is an integration convention, not an additional
equation of motion.
\leandef{MonadFormal.Temporal.Variational.Interval.unitMeasure,MonadFormal.Temporal.Variational.Interval.Velocity,MonadFormal.Temporal.Variational.Interval.primitive,MonadFormal.Temporal.Variational.Interval.primitiveCLM,MonadFormal.Temporal.Variational.Interval.endpoint,MonadFormal.Temporal.Variational.Interval.variationSpace,MonadFormal.Temporal.Variational.Interval.path}

\begin{lemma}[Bounded primitives and endpoint variations]
\label[lemma]{lem:variational-paths}
The map $P:V\to C_b(\mathbb R)$ is bounded linear, with
$\lVert Pf\rVert_\infty\le\lVert f\rVert_2$.
The subspace $H$ is closed. Each $x_f$ is absolutely continuous on
$[0,1]$, has endpoints $x_f(0)=a$, $x_f(1)=b$, and satisfies
$x_f'=b-a+f$ almost everywhere on that interval. Moreover,
\[
 \int f\,d\mu=0,\qquad
 \int (b-a+f)^2\,d\mu=(b-a)^2+\lVert f\rVert_2^2
 \quad(f\in H).
\]
\end{lemma}
\begin{proof}
Since $\mu(\mathbb R)=1$, Cauchy--Schwarz bounds each oriented
primitive by $\lVert f\rVert_2$. Integration is linear, and an
integrable function has a continuous primitive. Evaluation at $1$
is bounded, so its composite with $P$ has closed kernel $H$.
Extend $f$ by zero outside $[0,1]$. Ordinary integration of this
extension equals $Pf$. The fundamental theorem for Lebesgue
integrals gives absolute continuity and derivative $f$ almost
everywhere. Add the affine path to obtain the endpoint and derivative
claims. Finally, $(Pf)(1)=\int f\,d\mu=0$ removes the cross term
in the expansion of $(b-a+f)^2$.
\leanproof{MonadFormal.Temporal.Variational.Interval.velocity_integrable,MonadFormal.Temporal.Variational.Interval.primitive_continuous,MonadFormal.Temporal.Variational.Interval.primitive_bound,MonadFormal.Temporal.Variational.Interval.primitive_add,MonadFormal.Temporal.Variational.Interval.primitive_smul,MonadFormal.Temporal.Variational.Interval.primitiveCLM_apply,MonadFormal.Temporal.Variational.Interval.variation_initial,MonadFormal.Temporal.Variational.Interval.variation_terminal,MonadFormal.Temporal.Variational.Interval.extendedVelocity_integrable,MonadFormal.Temporal.Variational.Interval.primitive_eq_integral,MonadFormal.Temporal.Variational.Interval.primitive_absolutelyContinuous,MonadFormal.Temporal.Variational.Interval.primitive_hasDerivAt_ae,MonadFormal.Temporal.Variational.Interval.path_endpoints,MonadFormal.Temporal.Variational.Interval.path_hasDerivAt_ae,MonadFormal.Temporal.Variational.Interval.path_absolutelyContinuous,MonadFormal.Temporal.Variational.Interval.primitive_one_eq_integral,MonadFormal.Temporal.Variational.Interval.variation_integral_zero,MonadFormal.Temporal.Variational.Interval.integral_sq_eq_norm_sq,MonadFormal.Temporal.Variational.Interval.kinetic_integral}
\end{proof}

Fix $0\le\delta\le1$ and let $\nu_\delta$ be Lebesgue measure
restricted to $[0,1-\delta]$. Define
\[
 T_\delta f(t)=Pf(t+\delta)-Pf(t),\qquad
 b_\delta(t)=(b-a)\delta.
\]
Here $T_\delta:H\to L^2(\nu_\delta;\mathbb R)$ is bounded:
it composes $P$, translation minus the identity, and inclusion of
bounded continuous functions into $L^2$ of a finite measure.
\leandef{MonadFormal.Temporal.Variational.Interval.shiftMeasure,MonadFormal.Temporal.Variational.Interval.differenceCLM,MonadFormal.Temporal.Variational.Interval.shiftCLM,MonadFormal.Temporal.Variational.Interval.affineOffset,MonadFormal.Temporal.Variational.Interval.intervalEnergy}

\begin{theorem}[Nonlocal minimization in velocity coordinates]
\label{thm:interval-minimum}
For $\lambda\ge0$, the energy
\[
 J(x_f)=\frac12\int (b-a+f)^2\,d\mu
  +\frac\lambda2\int\bigl(x_f(t+\delta)-x_f(t)\bigr)^2\,d\nu_\delta(t)
\]
has a unique minimizing velocity class $f_*\in H$. It satisfies
\[
 \langle f_*,h\rangle_{L^2(\mu)}
 +\lambda\langle T_\delta f_*+b_\delta,T_\delta h\rangle_{L^2(\nu_\delta)}=0
 \quad(h\in H).
\]
\end{theorem}
\begin{proof}
Almost everywhere,
$T_\delta f+b_\delta=x_f(\,\cdot+\delta)-x_f$.
The preceding lemma therefore gives
\[
 J(x_f)=\frac{(b-a)^2}{2}
       +\frac{\lVert f\rVert_2^2}{2}
       +\frac\lambda2\lVert T_\delta f+b_\delta\rVert_2^2.
\]
Apply \Cref{thm:variational-minimum} to the Hilbert space $H$,
operator $T_\delta$ and offset $b_\delta$. The constant term does
not affect minimization.
\leanproof{MonadFormal.Temporal.Variational.Interval.shiftCLM_apply_ae,MonadFormal.Temporal.Variational.Interval.shifted_path_ae,MonadFormal.Temporal.Variational.Interval.intervalEnergy_eq_integrals,MonadFormal.Temporal.Variational.Interval.exists_unique_interval_minimizer,MonadFormal.Temporal.Variational.Interval.exists_interval_stationary_minimizer}
\end{proof}

% Lean module: MonadFormal.Temporal.Variational.Paths
\subsubsection{The full admissible path class}
\label{sec:variational-admissible}

Let $\mathcal A_{a,b}$ consist of absolutely continuous functions
$x:[0,1]\to\mathbb R$ with $x'\in L^2(\mu)$ and endpoints
$x(0)=a$, $x(1)=b$. These are the absolutely continuous
representatives of the endpoint-constrained one-dimensional $H^1$
class. Define
\[
 J(x)=\frac12\int (x')^2\,d\mu
      +\frac\lambda2\int\bigl(x(t+\delta)-x(t)\bigr)^2\,d\nu_\delta(t).
\]
Only values on $[0,1]$ enter this functional. In particular, extending
a path outside that interval changes neither its derivative almost
everywhere on the interval nor its shifted difference on
$[0,1-\delta]$.
\leandef{MonadFormal.Temporal.Variational.Interval.Admissible,MonadFormal.Temporal.Variational.Interval.pathEnergy}

\begin{theorem}[Exact velocity coordinates]
\label{thm:variational-coordinates}
The map $f\mapsto x_f|_{[0,1]}$ is a bijection from $H$ onto
$\mathcal A_{a,b}$. Its inverse is the $L^2$ class
$f=x'-(b-a)$. The path energy equals the coordinate energy in
\Cref{thm:interval-minimum}.
\end{theorem}
\begin{proof}
The construction gives an admissible path by
\Cref{lem:variational-paths}. Conversely, if $x\in\mathcal A_{a,b}$,
then $f=x'-(b-a)$ belongs to $L^2(\mu)$. Absolute continuity and
the fundamental theorem of calculus give, for $0\le t\le1$,
\[
 (Pf)(t)=x(t)-a-(b-a)t.
\]
At $t=1$ this is zero, so $f\in H$ and $x=x_f$ on $[0,1]$.
If two constructed paths agree there, their derivatives agree at
every interior point where both derivatives exist. Their velocity
classes therefore agree almost everywhere, giving uniqueness.
Finally, $x_f'=b-a+f$ almost everywhere identifies the kinetic
integral, while the shifted-path integral is unchanged.
\leanproof{MonadFormal.Temporal.Variational.Interval.unit_integral_eq,MonadFormal.Temporal.Variational.Interval.ae_mem_interior,MonadFormal.Temporal.Variational.Interval.deriv_eq_ae_of_eqOn,MonadFormal.Temporal.Variational.Interval.path_injective,MonadFormal.Temporal.Variational.Interval.path_admissible,MonadFormal.Temporal.Variational.Interval.primitive_centered_derivative,MonadFormal.Temporal.Variational.Interval.exists_unique_velocity,MonadFormal.Temporal.Variational.Interval.pathEnergy_path,MonadFormal.Temporal.Variational.Interval.pathEnergy_congr}
\end{proof}

\begin{theorem}[Unique admissible minimizing path]
\label{thm:variational-path-minimum}
For $0\le\delta\le1$ and $\lambda\ge0$, $J$ has a unique
minimizer in $\mathcal A_{a,b}$.
\end{theorem}
\begin{proof}
Transfer the minimizing velocity from \Cref{thm:interval-minimum}
through the energy-preserving bijection of
\Cref{thm:variational-coordinates}. Every admissible competitor
has a velocity coordinate, and every minimizing path must have
the unique minimizing coordinate. Thus uniqueness is equality
on $[0,1]$, with no assertion about an extension outside it.
\leanproof{MonadFormal.Temporal.Variational.Interval.exists_unique_path_minimizer}
\end{proof}

\input{MonadFormal/Temporal/Variational/Shift}
\input{MonadFormal/Temporal/Variational/Regularity}
\input{MonadFormal/Temporal/Variational/Euler}



